% Part: first-order-logic
% Chapter: model-theory
% Section: substructures

\documentclass[../../../include/open-logic-section]{subfiles}

\begin{document}

\olfileid[fr]{mod}{bas}{sub}
\olsection{Sous-\printtoken{p}{structure}}

Le !!{domain} d'!!a{structure}~$\Struct{M}$ peut être un sous-ensemble
de celui d'une autre !!{structure}~$\Struct{M'}$. Mais il ne suffit
évidemment pas que $\Domain{M} \subseteq \Domain{M'}$ pour considérer
$\Struct{M}$ comme une « partie » de $\Struct{M'}$ : il faut aussi que
$\Struct{M}$ et $\Struct{M'}$ s'accordent sur l'interprétation des
symboles du langage, du moins sur leur partie commune~$\Domain{M}$.

\begin{defn}
\ollabel{defn:substructure}
Soient $\Struct M$ et $\Struct M'$ des !!{structure}s pour un même
langage~$\Lang L$. On dit que $\Struct M$ est une
\emph{sous-!!{structure}} de $\Struct M'$, et que $\Struct M'$ est une
\emph{extension} de $\Struct M$, ce que l'on note
$\Struct M \substruct \Struct M'$, si et seulement si :
\begin{enumerate}
\item $\Domain{M} \subseteq \Domain{M'}$;
\item pour chaque constante $c \in \Lang L$, $\Assign{c}{M} =
    \Assign{c}{M'}$;
\item pour chaque !!{function} $n$-aire $f \in \Lang L$,
  $\Assign{f}{M}(a_1, \dots, a_n) = \Assign{f}{M'}(a_1, \dots, a_n)$
  pour tous $a_1$, \dots, $a_n \in \Domain{M}$;
\item pour chaque !!{predicate} $n$-aire $R \in \Lang L$,
  $\langle a_1, \dots, a_n\rangle \in \Assign{R}{M}$ si et seulement
  si $\langle a_1, \dots, a_n\rangle \in \Assign{R}{M'}$, pour tous
  $a_1$, \dots, $a_n \in \Domain{M}$.
\end{enumerate}
\end{defn}

\begin{rem}
\ollabel{rem:substructure}
Si le langage ne contient ni constante ni !!{function}, tout
sous-ensemble \emph{non vide} $N \subseteq \Domain{M}$ détermine une
sous-!!{structure}~$\Struct{N}$ de $\Struct M$, de !!{domain}
$\Domain{N} = N$, en posant $\Assign{R}{N} = \Assign{R}{M} \cap N^n$.\footnote{Précision éditoriale : la source autorise tout sous-ensemble ; nous excluons le sous-ensemble vide, puisque les domaines des structures sont non vides.}
\end{rem}

% Faut-il démontrer quelque chose ici ? Ajouter des exemples ?

\end{document}
